\documentclass[12pt]{article}
\usepackage{graphicx}

% Font size commands for consistency
\newcommand{\MainHeading}[1]{{\LARGE \textbf{#1}}}
\newcommand{\SubHeading}[1]{{\large \textbf{#1}}}

\begin{document}

\begin{titlepage}
\noindent
\includegraphics[width=2.5cm]{logo.png} 
\hspace{0.5cm}
\raisebox{0.9cm}{\MainHeading{Namal University, Mianwali}}

\vspace{0.3cm}

\begin{center}
{\Large \textbf{Project Proposal Document}}\\[1cm]

\SubHeading{Project:} Point-of-Sale System\\[0.3cm]
\SubHeading{Course:} CSC-225 – Software Engineering\\[0.3cm]
\SubHeading{Department:} Computer Science\\[0.3cm]
\SubHeading{Semester:} Fall 2025\\[1.5cm]

\textbf{Submitted To:}\\
Ms. Asiya Batool\\
Lecturer, Department of Computer Science\\[1.3cm]

\textbf{Submitted By:}\\
Muhammad Shoaib \hspace{0.5cm} Roll No: \underline{\hspace{3cm}}\\
Allah Ditta \hspace{1.9cm} Roll No: \underline{\hspace{3cm}}\\
Zainab Bibi \hspace{1.9cm} Roll No: \underline{\hspace{3cm}}\\[0.5cm]

\vfill
\textbf{Submission Date: 9 Nov 2015} \\
\end{center}

\end{titlepage}


\begin{center}
\pagestyle{empty}
\SubHeading{Requirement Provider Agreement}
\end{center}
\noindent
\SubHeading{Project Title:} Point-of-Sale System\\
\SubHeading{Course:} CSC-225 – Software Engineering\\
\SubHeading{Department:} Computer Science, Namal University

\vspace{1cm}

\noindent
This agreement is made between the following parties.

\vspace{0.3cm}

\noindent
\SubHeading{Student Team:}\\
Muhammad Shoaib \hspace{1cm}Reg No: \underline{\hspace{4cm}}\\
Allah Ditta \hspace{2.45cm}Reg No: \underline{\hspace{4cm}}\\
Zainab Bibi \hspace{2.4cm}Reg No: \underline{\hspace{4cm}}\\

\vspace{0.5cm}

\noindent
\SubHeading{Requirement Provider (RP):}\\
Dr Muzamil Ahmed\\
Assistant Professor, Department of Computer Science, Namal University

\vspace{1cm}

\noindent
\SubHeading{Agreement Purpose:}\\
The student team and the Requirement Provider agree to work together for the Software Engineering Project. Dr Muzamil will provide requirements, feedback, and guidance. The team will do regular meetings and share progress updates.

\vspace{1cm}

\noindent
\SubHeading{Signatures:}\\
Muhammad Shoaib (Team Leader)\hspace{0.5cm}Signature: \underline{\hspace{4cm}}\\[0.1cm]
Dr Muzamil Ahmed (RP) \hspace{2.05cm}Signature: \underline{\hspace{4cm}}\\[0.5cm]
Date: \hspace{7.5cm} \underline{\hspace{4cm}}

\newpage
\tableofcontents
\newpage
\section{Introduction}
Restaurants, such as pizza shops, face difficulties in managing orders, sales, and inventory when operating multiple branches.which is where a Point-of-Sale (POS) system can help.A Point-of-Sale (POS) system makes the running of restaurants much easier.It helps in taking orders, handling payments, and keeping everything organized even in busy hours.Shows menu and stock in one place, so even if different branches have different menus, the system  shows the correct items for each branch.It also helps to track which dishes sell well, alerts staff when ingredients are running low, and ensures that prices are consistent across all branches,reducing mistakes, and making service faster and smoother.
\section{Problem Statement}
Mostly, the pizza shop manually handles all their managements.which often creates several problems.Orders are usually written by hand, leading to mistakes such as wrong items, missing custom requests or miscommunication between waiter and kitchen staff, a worker may take an order for a pizza that exists at another branch but not at their own if the restaurant has different branches. Prices or ingredients might also different and without a proper system its hard for staff to remember all those details.Manual billing and order tracking take more time,especially during busy hours,leading delay in service.Without a POS system its also difficult to get accurate sales reports, identify the best selling items,or monitoring staff performance.Customer may even miss out better deals or discount is their order is similar in value to a deal but is entered separately at a higher price.

\end{document}
